\section{Part 1: Vertical variation of wind speed}
\label{sec:part1}

\subsection{Met Mast and Location}
\textbf{Task:} Describe the met mast and its location. How do the surroundings of the mast look in terms of terrain roughness? How far is the mast from the turbine and in which direction? Why?

The measurements are from DTU Risø Campus using a met mast with a total height of 70 m. The mast is around 120 m, away from the Vestas V52 wind turbine.
The surrounding terrain has a low roughness to the west  (due to water from the Roskilde Fjord) and a higher roughness in the other directions since it consists of open fields and the campus buildings.
To assure the free flow, the mast is posi
% Hier Text einfügen

\subsection{Data Collection}
\textbf{Task:} Describe the data collected. What has been measured, what units, what sample rate, where at the mast and with what instruments?

\subsection{Raw Time Series Plot}
\textbf{Task:} Plot the raw time series of the wind speed at the different heights. What time scales do you see? Do they differ for the different heights?

\begin{figure}[H]
    \centering
    % \includegraphics[width=0.8\textwidth]{figures/Q3_plot.png}
    \caption{Raw wind speed time series at various heights (10th Feb 2026).}
\end{figure}

\subsection{Logarithmic Wind Profile}
\textbf{Task:} Plot the average wind speed for each height against height $z$. Can the profile be described by a logarithmic wind profile?

\subsection{Roughness Length and Friction Velocity}
\textbf{Task:} Estimate the roughness length $z_0$ and the friction velocity $u_*$ from the observed wind profile.

\subsection{Stability Impact}
\textbf{Task:} Discuss the validity of the logarithmic profile under the observed stability conditions.

\subsection{Heat Flux and Stability}
\textbf{Task:} Estimate the friction velocity and the heat flux from the sonic anemometer data.

\subsection{Weather Conditions}
\textbf{Task:} How was the weather at the day of visit? Does the observed weather support your finding of question 7?


The day we visited Risø was windy from the eastern direction with some snowflakes. Generally, the site experiences western winds from the fjord, but that was not the case that day.

\textcolor{red}{does it support finding of Q7?}
